% back/app.tex
\chapter{附錄}

\begingroup
% 讓 A/B/C 用字母當小節序號（顯示成 A、B、C）
\renewcommand{\thesection}{\Alph{section}}
\renewcommand{\thesubsection}{\thesection.\arabic{subsection}}
\setcounter{section}{0}

\section{文字的數學結構}
爲了防止在正文中過多的技術細節，我們將數學化的理論模型放到附錄來解決。

\section{與文字相關的符號系統}
有許多符號系統具有與文字類似的性質，但卻與我們所探討的內容有所差距，因此放到附錄來作一個綜述。

\section{Unicode對漢字的處理}
漢字在Unicode中必然具有特殊的地位，也涉及相當多複雜的部分。在正文中我們已經對Unicode作過宏觀的介紹，但漢字其實很適合單獨拿出來處理，因此放在附錄中。

\section{有關文字史的補充}
正文中無法涵蓋所有的文字，也就不可能介紹所有文字的歷史，而且按類型學展開的介紹往往很難有歷史的系統性。我們在附錄中對整個文字史進行串講，希望能夠
建立對文字史的整體把握。在正文中提到過的文字，我們會簡單帶過。

\section{國際音標簡介}
正文將國際音標默認爲背景知識，但對於不熟悉其的人，仍然有必要進行一個簡短的介紹。

\endgroup
